% !TEX program = pdflatex
% FRE6883 Recitation 3 -- Pointers, Function Pointers, and Generic CRR
\documentclass[10pt,aspectratio=169]{beamer}

\usepackage[utf8]{inputenc}
\usepackage[T1]{fontenc}
\usepackage[english]{babel}
\usepackage{lmodern}
\usepackage{microtype}
\DisableLigatures{encoding=*,family=tt*}
\usepackage{amsmath,amssymb,bm,mathtools}
\usepackage{booktabs}
\usepackage{multicol}
\usepackage{changepage}
\usepackage{graphicx}
\usepackage{tikz}
\usetikzlibrary{arrows.meta,positioning,calc,decorations.pathreplacing}
\usepackage{pgfplots}
\pgfplotsset{compat=1.18}

\usetheme[numbering=fraction,progressbar=frametitle]{metropolis}
\setbeamertemplate{frame numbering}[fraction]
\setbeamertemplate{navigation symbols}{}

\definecolor{BootcampBlueDark}{HTML}{1A3C68}
\definecolor{BootcampBlue}{RGB}{31,110,178}
\definecolor{AccentGold}{HTML}{DAA520}
\definecolor{SoftGray}{RGB}{245,246,248}
\definecolor{TextBlack}{HTML}{000000}
\definecolor{Crimson}{HTML}{990000}
\definecolor{MatrixGreen}{HTML}{228B22}
\definecolor{MatrixRed}{HTML}{B22222}

\setbeamercolor{normal text}{fg=TextBlack,bg=white}
\setbeamercolor{title}{fg=TextBlack}
\setbeamercolor{structure}{fg=BootcampBlueDark}
\setbeamercolor{progress bar}{fg=BootcampBlueDark,bg=gray!30}
\setbeamercolor{frametitle}{bg=BootcampBlueDark,fg=white}
\setbeamerfont{frametitle}{series=\bfseries,size=\large}
\setbeamertemplate{frametitle}{
  \nointerlineskip
  \begin{beamercolorbox}[wd=\paperwidth,ht=3.2ex,dp=1.4ex,leftskip=1em,rightskip=1em]{frametitle}
    \usebeamerfont{frametitle}\insertframetitle
  \end{beamercolorbox}
}

\setbeamercolor{block title example}{bg=BootcampBlueDark!90!black,fg=white}
\setbeamercolor{block body example}{bg=BootcampBlueDark!5!white,fg=black}
\setbeamercolor{block title proposition}{bg=MatrixGreen!75!black,fg=white}
\setbeamercolor{block body proposition}{bg=MatrixGreen!7!white,fg=black}
\setbeamercolor{block title illustration}{bg=AccentGold!85!black,fg=white}
\setbeamercolor{block body illustration}{bg=AccentGold!10!white,fg=black}
\setbeamercolor{block title proof}{bg=gray!70!black,fg=white}
\setbeamercolor{block body proof}{bg=SoftGray,fg=black}
\setbeamertemplate{blocks}[rounded][shadow=false]

\newenvironment{definitionblock}[1]{%
  \par\vspace{0.3em}%
  \begin{beamercolorbox}[wd=\textwidth,sep=0.35em,leftskip=0.3em,rightskip=0.6em,rounded=false,shadow=false]{block title example}%
  \raisebox{0.15em}{\usebeamerfont{block title example}#1}%
  \end{beamercolorbox}\vspace{-0.4em}%
  \begin{beamercolorbox}[wd=\textwidth,sep=0.55em,leftskip=0.6em,rightskip=0.6em,rounded=false,shadow=false]{block body example}%
  \usebeamerfont{block body example}%
}{\end{beamercolorbox}\par\vspace{0.45em}}

\newenvironment{propositionblock}[1]{%
  \par\vspace{0.3em}%
  \begin{beamercolorbox}[wd=\textwidth,sep=0.35em,leftskip=0.3em,rightskip=0.6em,rounded=false,shadow=false]{block title proposition}%
  \raisebox{0.15em}{\usebeamerfont{block title example}#1}%
  \end{beamercolorbox}\vspace{-0.4em}%
  \begin{beamercolorbox}[wd=\textwidth,sep=0.55em,leftskip=0.6em,rightskip=0.6em,rounded=false,shadow=false]{block body proposition}%
  \usebeamerfont{block body example}%
}{\end{beamercolorbox}\par\vspace{0.45em}}

\newenvironment{illustrationblock}[1]{%
  \par\vspace{0.3em}%
  \begin{beamercolorbox}[wd=\textwidth,sep=0.35em,leftskip=0.3em,rightskip=0.6em,rounded=false,shadow=false]{block title illustration}%
  \raisebox{0.15em}{\usebeamerfont{block title example}#1}%
  \end{beamercolorbox}\vspace{-0.4em}%
  \begin{beamercolorbox}[wd=\textwidth,sep=0.55em,leftskip=0.6em,rightskip=0.6em,rounded=false,shadow=false]{block body illustration}%
  \usebeamerfont{block body example}%
}{\end{beamercolorbox}\par\vspace{0.45em}}

\newenvironment{proofblock}[1]{%
  \par\vspace{0.3em}%
  \begin{beamercolorbox}[wd=\textwidth,sep=0.35em,leftskip=0.3em,rightskip=0.6em,rounded=false,shadow=false]{block title proof}%
  \raisebox{0.15em}{\usebeamerfont{block title example}#1}%
  \end{beamercolorbox}\vspace{-0.4em}%
  \begin{beamercolorbox}[wd=\textwidth,sep=0.55em,leftskip=0.6em,rightskip=0.6em,rounded=false,shadow=false]{block body proof}%
  \usebeamerfont{block body example}%
}{\end{beamercolorbox}\par\vspace{0.45em}}

\usepackage{listings}
\usepackage{textcomp}
\lstdefinestyle{coursecpp}{
  language=C++,basicstyle=\ttfamily\fontsize{9.1}{10.4}\selectfont,
  keywordstyle=\color{blue},commentstyle=\color{MatrixGreen},stringstyle=\color{Crimson},
  showstringspaces=false,columns=fullflexible,keepspaces=true,tabsize=4,
  breaklines=false,frame=none,aboveskip=0.25em,belowskip=0.25em,upquote=true
}
\lstdefinestyle{solutioncpp}{style=coursecpp,basicstyle=\ttfamily\fontsize{8.3}{9.4}\selectfont}
\lstdefinestyle{tinycpp}{style=coursecpp,basicstyle=\ttfamily\fontsize{7.5}{8.5}\selectfont}
\lstset{style=coursecpp}
\setbeamercovered{invisible}
\setlength{\parskip}{0.1em}

\newcommand{\problemmeta}[2]{{\color{BootcampBlue}\small\textbf{GUIDED CODE}\quad #1\hfill #2}\par\smallskip\small}
\newcommand{\explain}[2]{{\color{BootcampBlueDark}\textbf{#1}}\par #2\par\medskip}

\setbeamertemplate{frame footer}{\scriptsize FRE6883 | Recitation 3 | Pointers and Generic CRR}
\title{Pointers, Function Pointers, and a Generic CRR Pricer}
\author{}
\date{}

% 45 core frames for approximately 100 minutes, followed by reserve brainteasers.
% Source baseline: Lecture 3, Topic 3, slides 2--47.
\begin{document}

% FRAME 1
\begin{frame}[plain]
\vfill
{\usebeamerfont{title}\inserttitle\par}
\medskip
{\color{BootcampBlueDark}\rule{0.95\textwidth}{0.8pt}}
\bigskip
{\Large FRE6883 -- Recitation 3\par}
\medskip
Trace the address, protect the lifetime, choose the payoff
\vfill
\end{frame}

% FRAME 2
\begin{frame}{One address lets us change data and behavior}
\begin{columns}[c,onlytextwidth]
\begin{column}{0.46\textwidth}
\centering
Variable address $\rightarrow$ pointer\par
$\downarrow$\par
Arrays and pointer arithmetic\par
$\downarrow$\par
Dynamic memory and ownership\par
$\downarrow$\par
Function address $\rightarrow$ payoff strategy
\end{column}
\begin{column}{0.50\textwidth}
\begin{propositionblock}{By the end}
You can trace a pointer expression, identify unsafe lifetime behavior, and reuse one CRR engine for calls, puts, and digital payoffs.
\end{propositionblock}
\medskip
Every code activity follows an example from Lecture 3, then asks for one prediction or safety check.
\end{column}
\end{columns}
\end{frame}

% FRAME 3
\begin{frame}[fragile]{A pointer stores where a value lives}
\begin{columns}[T,onlytextwidth]
\begin{column}{0.56\textwidth}
\begin{lstlisting}
double dPrice = 21.68;
double* dPtr = &dPrice;

std::cout << dPtr;   // an address
std::cout << *dPtr;  // 21.68
\end{lstlisting}
\end{column}
\begin{column}{0.40\textwidth}
\small
\explain{\texttt{\&dPrice}}{The address of the variable.}
\explain{\texttt{double* dPtr}}{A pointer that may point to a \texttt{double}.}
\explain{\texttt{*dPtr}}{Dereference the pointer: access the value at that address.}
\end{column}
\end{columns}
\begin{illustrationblock}{Keep the two meanings of \texttt{*} separate}
In a declaration, \texttt{*} introduces a pointer. In an expression, \texttt{*} follows the stored address.
\end{illustrationblock}
\end{frame}

% FRAME 4
\begin{frame}[fragile]{01 | Complete the address and dereference operations -- Problem}
\problemmeta{Pointer basics}{Predict before compiling}
\begin{lstlisting}
double dPrice = 21.68;
double* dPtr = /* address of dPrice */;

std::cout << /* print the address */ << std::endl;
std::cout << /* print 21.68 */ << std::endl;

/* change dPrice through dPtr */ = 24.50;
std::cout << dPrice << std::endl;
\end{lstlisting}
\begin{definitionblock}{Expected final output}
The last line prints \texttt{24.5}. Explain why no assignment to the name \texttt{dPrice} was needed.
\end{definitionblock}
\end{frame}

% FRAME 5
\begin{frame}[fragile]{01 | Complete the address and dereference operations -- Solution}
\begin{columns}[T,onlytextwidth]
\begin{column}{0.61\textwidth}
\begin{lstlisting}[style=solutioncpp]
double dPrice = 21.68;
double* dPtr = &dPrice;

std::cout << dPtr << std::endl;
std::cout << *dPtr << std::endl;

*dPtr = 24.50;
std::cout << dPrice << std::endl;
\end{lstlisting}
\end{column}
\begin{column}{0.35\textwidth}
\small
\explain{Alias}{\texttt{dPtr} and \texttt{dPrice} give two access paths to the same object.}
\explain{Type check}{A \texttt{double*} must point to a \texttt{double}.}
\explain{Address output}{The hexadecimal-looking address is machine- and run-dependent.}
\end{column}
\end{columns}
\end{frame}

% FRAME 6
\begin{frame}[fragile]{02 | Trace two aliases to one price -- Problem}
\problemmeta{Aliasing}{Fill the three outputs}
\begin{lstlisting}
double price = 100.0;
double* p = &price;
double* q = p;

*p += 5.0;
std::cout << price << " ";      // _____

*q *= 1.10;
std::cout << *p << " ";         // _____
std::cout << (p == q) << "\n"; // _____
\end{lstlisting}
\begin{illustrationblock}{Question}
Are there one, two, or three \texttt{double} objects in this example?
\end{illustrationblock}
\end{frame}

% FRAME 7
\begin{frame}[fragile]{02 | Trace two aliases to one price -- Solution}
\begin{lstlisting}[style=solutioncpp]
*p += 5.0;   // price is 105.0
*q *= 1.10;  // the same price is 115.5

// Output: 105 115.5 1
\end{lstlisting}
\begin{columns}[T,onlytextwidth]
\begin{column}{0.48\textwidth}
\begin{definitionblock}{One object}
There is one \texttt{double} and two pointer variables. Both pointers store the same address.
\end{definitionblock}
\end{column}
\begin{column}{0.48\textwidth}
\begin{propositionblock}{Debugging consequence}
A write through either alias changes what every other alias observes.
\end{propositionblock}
\end{column}
\end{columns}
\end{frame}

% FRAME 8
\begin{frame}[fragile]{An array name identifies its first element}
\begin{columns}[T,onlytextwidth]
\begin{column}{0.57\textwidth}
\begin{lstlisting}
int values[] = {2, 4, 6};
int* ptr = values;

values[0]     // 2
*values       // 2
ptr[1]        // 4
*(ptr + 2)    // 6
\end{lstlisting}
\end{column}
\begin{column}{0.39\textwidth}
\small
\explain{Key equivalence}{\texttt{values[i]} means \texttt{*(values + i)}.}
\explain{Scaled movement}{Adding one moves by one element, not one byte.}
\explain{No guardrail}{Raw array access performs no bounds checking.}
\end{column}
\end{columns}
\end{frame}

% FRAME 9
\begin{frame}[fragile]{03 | Write the same array access four ways -- Problem}
\problemmeta{Arrays and pointers}{Make every line assign 12 to index 2}
\begin{lstlisting}
int values[] = {2, 4, 6};
int* ptr = values;

values[/* ? */] = 12;
ptr[/* ? */] = 12;
*(values + /* ? */) = 12;
*(ptr + /* ? */) = 12;
\end{lstlisting}
\begin{definitionblock}{Constraint}
Do not change the operators. Only complete the four indices or offsets.
\end{definitionblock}
\end{frame}

% FRAME 10
\begin{frame}[fragile]{03 | Write the same array access four ways -- Solution}
\begin{lstlisting}[style=solutioncpp]
values[2] = 12;
ptr[2] = 12;
*(values + 2) = 12;
*(ptr + 2) = 12;
\end{lstlisting}
\begin{propositionblock}{The compiler sees the same location}
All four expressions target the third \texttt{int}. Choose array notation for readability; understand pointer notation so you can trace lower-level code.
\end{propositionblock}
\begin{illustrationblock}{Boundary check}
\texttt{values[3]} compiles, but it accesses outside this three-element array: the behavior is undefined.
\end{illustrationblock}
\end{frame}

% FRAME 11
\begin{frame}[fragile]{Parentheses decide whether we move or add}
\begin{columns}[T,onlytextwidth]
\begin{column}{0.54\textwidth}
\begin{lstlisting}
int values[] = {2, 4, 6};
int* ptr = values;

*(ptr + 1)  // 4
*ptr + 1    // 3
\end{lstlisting}
\end{column}
\begin{column}{0.42\textwidth}
\begin{definitionblock}{Move, then dereference}
\texttt{*(ptr + 1)} selects the next element.
\end{definitionblock}
\begin{definitionblock}{Dereference, then add}
\texttt{*ptr + 1} reads the first value and adds one to it.
\end{definitionblock}
\end{column}
\end{columns}
\end{frame}

% FRAME 12
\begin{frame}[fragile]{04 | Trace pointer arithmetic without running it -- Problem}
\problemmeta{Movement and distance}{Complete the outputs}
\begin{lstlisting}
int values[] = {2, 4, 6};
int* ptr = values;

std::cout << *(++ptr) << " "; // _____
ptr += 1;
std::cout << *ptr << " ";     // _____
std::cout << ptr - values;     // _____
\end{lstlisting}
\begin{illustrationblock}{Invariant}
At every step, state the index at which \texttt{ptr} points before stating the value.
\end{illustrationblock}
\end{frame}

% FRAME 13
\begin{frame}[fragile]{04 | Trace pointer arithmetic without running it -- Solution}
\begin{lstlisting}[style=solutioncpp]
std::cout << *(++ptr) << " "; // ptr -> values[1], prints 4
ptr += 1;                     // ptr -> values[2]
std::cout << *ptr << " ";     // prints 6
std::cout << ptr - values;    // prints 2
\end{lstlisting}
\begin{definitionblock}{Pointer subtraction}
For two pointers into the same array, subtraction reports the number of elements between them.
\end{definitionblock}
\begin{propositionblock}{Output}
\texttt{4 6 2}
\end{propositionblock}
\end{frame}

% FRAME 14
\begin{frame}[fragile]{Comparing pointers is different from comparing values}
\begin{lstlisting}
if (p == q)      // same address?
if (*p == *q)    // same pointed-to value?
\end{lstlisting}
\begin{columns}[T,onlytextwidth]
\begin{column}{0.48\textwidth}
\begin{definitionblock}{Address equality}
\texttt{p == q} asks whether the two pointers identify the same object.
\end{definitionblock}
\end{column}
\begin{column}{0.48\textwidth}
\begin{definitionblock}{Value equality}
\texttt{*p == *q} dereferences both pointers and compares the two values.
\end{definitionblock}
\end{column}
\end{columns}
\end{frame}

% FRAME 15
\begin{frame}[fragile]{05 | Predict address equality and value equality -- Problem}
\problemmeta{Pointer comparison}{Fill both booleans}
\begin{lstlisting}
int x = 7;
int y = 7;
int* p = &x;
int* q = &y;

std::cout << (p == q) << " ";   // _____
std::cout << (*p == *q) << "\n"; // _____
\end{lstlisting}
\begin{illustrationblock}{Explain before revealing}
The integers begin with equal values, but they are distinct objects.
\end{illustrationblock}
\end{frame}

% FRAME 16
\begin{frame}[fragile]{05 | Predict address equality and value equality -- Solution}
\begin{lstlisting}[style=solutioncpp]
std::cout << (p == q) << " ";    // 0: distinct addresses
std::cout << (*p == *q) << "\n"; // 1: equal values
\end{lstlisting}
\begin{propositionblock}{Output}
\texttt{0 1}
\end{propositionblock}
\begin{illustrationblock}{Safety precondition}
Only dereference a pointer when it points to a live object. Comparing a pointer with \texttt{nullptr} does not dereference it.
\end{illustrationblock}
\end{frame}

% FRAME 17
\begin{frame}[fragile]{A pointer parameter can modify the caller's object}
\begin{columns}[T,onlytextwidth]
\begin{column}{0.55\textwidth}
\begin{lstlisting}
void getNum(int* ptr)
{
    std::cin >> *ptr;
}

int n = 1;
getNum(&n);
\end{lstlisting}
\end{column}
\begin{column}{0.41\textwidth}
\small
\explain{At the call}{\texttt{\&n} passes the address of the caller's integer.}
\explain{Inside}{\texttt{*ptr} accesses that integer.}
\explain{Effect}{The assignment changes \texttt{n}, just as a non-const reference parameter would.}
\end{column}
\end{columns}
\end{frame}

% FRAME 18
\begin{frame}[fragile]{06 | Complete call by pointer -- Problem}
\problemmeta{Lecture 3 \texttt{getNum}}{Make the caller observe 9}
\begin{lstlisting}
void getNum(int* ptr)
{
    std::cin >> /* pointed-to integer */;
}

int main()
{
    int number = 1;
    getNum(/* address of number */);
    std::cout << number << std::endl;
}
\end{lstlisting}
\begin{definitionblock}{Input}
If the user enters \texttt{9}, the program must print \texttt{9}.
\end{definitionblock}
\end{frame}

% FRAME 19
\begin{frame}[fragile]{06 | Complete call by pointer -- Solution}
\begin{lstlisting}[style=solutioncpp]
void getNum(int* ptr)
{
    std::cin >> *ptr;
}

int number = 1;
getNum(&number);
\end{lstlisting}
\begin{columns}[T,onlytextwidth]
\begin{column}{0.48\textwidth}
\begin{definitionblock}{Contract}
The pointer must not be \texttt{nullptr}; the function dereferences it immediately.
\end{definitionblock}
\end{column}
\begin{column}{0.48\textwidth}
\begin{propositionblock}{Reference comparison}
\texttt{void getNum(int\& number)} expresses the same required-object relationship more directly.
\end{propositionblock}
\end{column}
\end{columns}
\end{frame}

% FRAME 20
\begin{frame}[fragile]{07 | Swap two caller-owned integers -- Problem}
\problemmeta{Lecture 3 \texttt{MySwap}}{Complete three assignments}
\begin{lstlisting}
void MySwap(int* x, int* y)
{
    int temp = /* value pointed to by x */;
    /* object x */ = /* value pointed to by y */;
    /* object y */ = temp;
}

int a = 2, b = -3;
MySwap(/* two addresses */);
\end{lstlisting}
\begin{definitionblock}{Postcondition}
After the call, \texttt{a == -3} and \texttt{b == 2}.
\end{definitionblock}
\end{frame}

% FRAME 21
\begin{frame}[fragile]{07 | Swap two caller-owned integers -- Solution}
\begin{lstlisting}[style=solutioncpp]
void MySwap(int* x, int* y)
{
    int temp = *x;
    *x = *y;
    *y = temp;
}

int a = 2, b = -3;
MySwap(&a, &b);
\end{lstlisting}
\begin{propositionblock}{Reason about effects}
The pointer variables \texttt{x} and \texttt{y} are local copies of addresses, but dereferencing them reaches the caller's objects.
\end{propositionblock}
\end{frame}

% FRAME 22
\begin{frame}[fragile]{Read pointer declarations from right to left}
\small
\begin{center}
\renewcommand{\arraystretch}{1.25}
\begin{tabular}{@{}lll@{}}
\toprule
Declaration & Can move pointer? & Can change value through it? \\
\midrule
\texttt{double* p} & yes & yes \\
\texttt{const double* p} & yes & no \\
\texttt{double* const p = \&x} & no & yes \\
\texttt{const double* const p = \&x} & no & no \\
\bottomrule
\end{tabular}
\end{center}
\begin{definitionblock}{Two different promises}
\texttt{const} before the pointed-to type protects the value. \texttt{const} after \texttt{*} fixes the address stored in the pointer.
\end{definitionblock}
\end{frame}

% FRAME 23
\begin{frame}[fragile]{08 | Decide which operations compile -- Problem}
\problemmeta{Pointer to const versus const pointer}{Mark valid or invalid}
\begin{lstlisting}
double x = 10.0, y = 20.0;
const double* readOnly = &x;
double* const fixed = &x;

readOnly = &y;   // valid / invalid?
*readOnly = 7.0; // valid / invalid?
fixed = &y;      // valid / invalid?
*fixed = 7.0;    // valid / invalid?
\end{lstlisting}
\begin{illustrationblock}{Reading trick}
Ask two separate questions: may the address change, and may the pointed-to value change?
\end{illustrationblock}
\end{frame}

% FRAME 24
\begin{frame}[fragile]{08 | Decide which operations compile -- Solution}
\begin{lstlisting}[style=solutioncpp]
readOnly = &y;   // valid: pointer may move
*readOnly = 7.0; // invalid: value is protected
fixed = &y;      // invalid: pointer is fixed
*fixed = 7.0;    // valid: x becomes 7.0
\end{lstlisting}
\begin{columns}[T,onlytextwidth]
\begin{column}{0.48\textwidth}
\begin{definitionblock}{Useful parameter}
\texttt{const double* data} lets a function read an array while promising not to edit its elements.
\end{definitionblock}
\end{column}
\begin{column}{0.48\textwidth}
\begin{illustrationblock}{Important nuance}
The original variable may still be mutable through another non-const access path.
\end{illustrationblock}
\end{column}
\end{columns}
\end{frame}

% FRAME 25
\begin{frame}{Dynamic memory has an explicit lifetime}
\centering
\Large
\texttt{new / new[]} $\longrightarrow$ use the live object $\longrightarrow$ \texttt{delete / delete[]}
\normalsize
\vspace{1.3em}
\begin{columns}[T,onlytextwidth]
\begin{column}{0.47\textwidth}
\begin{definitionblock}{Scalar}
\texttt{double* p = new double;}\par
Release with \texttt{delete p;}.
\end{definitionblock}
\end{column}
\begin{column}{0.49\textwidth}
\begin{definitionblock}{Array}
\texttt{double* p = new double[n];}\par
Release with \texttt{delete[] p;}.
\end{definitionblock}
\end{column}
\end{columns}
\begin{propositionblock}{Ownership question}
Every successful allocation needs a clearly responsible owner that releases it exactly once.
\end{propositionblock}
\end{frame}

% FRAME 26
\begin{frame}[fragile]{09 | Allocate, use, and release one price -- Problem}
\problemmeta{Dynamic scalar}{Complete the lifetime}
\begin{lstlisting}
double* price = /* allocate one double */;
/* assign through the pointer */ = 21.68;
std::cout << /* pointed-to value */ << std::endl;

/* release scalar memory */;
price = /* safe empty state */;
\end{lstlisting}
\begin{definitionblock}{Expected output}
\texttt{21.68}
\end{definitionblock}
\end{frame}

% FRAME 27
\begin{frame}[fragile]{09 | Allocate, use, and release one price -- Solution}
\begin{lstlisting}[style=solutioncpp]
double* price = new double;
*price = 21.68;
std::cout << *price << std::endl;

delete price;
price = nullptr;
\end{lstlisting}
\begin{definitionblock}{After deletion}
The dynamic object no longer exists. Setting the pointer to \texttt{nullptr} prevents this variable from continuing to look usable.
\end{definitionblock}
\begin{illustrationblock}{It does not repair aliases}
Any other pointer that stored the old address is still dangling.
\end{illustrationblock}
\end{frame}

% FRAME 28
\begin{frame}[fragile]{10 | Build a dynamically sized payoff row -- Problem}
\problemmeta{Dynamic array}{Match every bracket}
\begin{lstlisting}
int N = 3;
double* payoff = new double[/* element count */];

for (int i = 0; i <= N; ++i)
{
    payoff[i] = static_cast<double>(i * i);
}

std::cout << payoff[3] << std::endl;
/* release array memory */;
payoff = nullptr;
\end{lstlisting}
\end{frame}

% FRAME 29
\begin{frame}[fragile]{10 | Build a dynamically sized payoff row -- Solution}
\begin{lstlisting}[style=solutioncpp]
int N = 3;
double* payoff = new double[N + 1];

for (int i = 0; i <= N; ++i)
{
    payoff[i] = static_cast<double>(i * i);
}

std::cout << payoff[3] << std::endl; // 9
delete[] payoff;
payoff = nullptr;
\end{lstlisting}
\begin{propositionblock}{Pairing rule}
\texttt{new[]} pairs with \texttt{delete[]}; scalar \texttt{new} pairs with scalar \texttt{delete}.
\end{propositionblock}
\end{frame}

% FRAME 30
\begin{frame}[fragile]{Leaks, dangling pointers, and invalid dereferences differ}
\small
\begin{center}
\renewcommand{\arraystretch}{1.2}
\begin{tabular}{@{}p{2.5cm}p{4.2cm}p{4.1cm}@{}}
\toprule
Failure & What happened & Prevention \\
\midrule
Memory leak & Last owning address is lost before \texttt{delete} & Preserve ownership; release once \\
Dangling pointer & Pointer still stores an address after object lifetime ends & Stop using it; set owned pointer to \texttt{nullptr} \\
Null dereference & Code evaluates \texttt{*p} when \texttt{p == nullptr} & Check the contract before dereferencing \\
Wrong return & Function returns the address of a local object & Return by value or return live dynamic/caller data \\
\bottomrule
\end{tabular}
\end{center}
\begin{illustrationblock}{Modern C++ note}
Lecture 3 teaches raw ownership directly. In production code, containers and smart pointers usually make the lifetime safer.
\end{illustrationblock}
\end{frame}

% FRAME 31
\begin{frame}[fragile]{11 | Classify four ways to return or modify a number -- Problem}
\problemmeta{Lifetime audit from Lecture 3}{Safe, unsafe, or null?}
\begin{lstlisting}[style=tinycpp]
int* getNum1() { int n = 10; return &n; }
int  getNum2() { int n = 10; return n + 1; }
void getNum3(int* p) { *p += 1; }
int* getNum4() {
    int* p = new int{11};
    delete p;
    p = nullptr;
    return p;
}
\end{lstlisting}
\begin{definitionblock}{Task}
For each function, state what it returns and whether the caller may safely dereference the result.
\end{definitionblock}
\end{frame}

% FRAME 32
\begin{frame}[fragile]{11 | Classify four ways to return or modify a number -- Solution}
\small
\begin{center}
\renewcommand{\arraystretch}{1.22}
\begin{tabular}{@{}p{1.7cm}p{3.0cm}p{5.8cm}@{}}
\toprule
Function & Result & Judgment \\
\midrule
\texttt{getNum1} & dangling address & unsafe: local \texttt{n} dies on return \\
\texttt{getNum2} & value \texttt{11} & safe: return by value \\
\texttt{getNum3} & no return value & safe when caller passes a valid address \\
\texttt{getNum4} & \texttt{nullptr} & cannot dereference; the allocated object was deleted \\
\bottomrule
\end{tabular}
\end{center}
\begin{propositionblock}{A safe returning pointer must identify live data}
That data must belong to the caller, have static lifetime, or remain dynamically allocated with a clear owner.
\end{propositionblock}
\end{frame}

% FRAME 33
\begin{frame}[fragile]{A function pointer stores which operation to call}
\begin{columns}[T,onlytextwidth]
\begin{column}{0.57\textwidth}
\begin{lstlisting}
int addition(int a, int b)
{
    return a + b;
}

int (*operation)(int, int) = addition;
int result = operation(1, 2);
\end{lstlisting}
\end{column}
\begin{column}{0.39\textwidth}
\small
\explain{Return type}{The pointed-to function returns \texttt{int}.}
\explain{Parameter list}{It accepts two \texttt{int} arguments.}
\explain{Behavior as data}{Assign another compatible function to change the operation without changing the calling code.}
\end{column}
\end{columns}
\end{frame}

% FRAME 34
\begin{frame}[fragile]{12 | Complete a function pointer declaration -- Problem}
\problemmeta{Lecture 3 addition example}{Match the exact signature}
\begin{lstlisting}
int addition(int a, int b)
{
    return a + b;
}

int (/* pointer syntax */)(int, int) = addition;
int result = /* invoke through pointer */(1, 2);
std::cout << result << std::endl;
\end{lstlisting}
\begin{illustrationblock}{Why the parentheses matter}
Without parentheses around \texttt{*operation}, the declaration describes a function returning a pointer, not a pointer to a function.
\end{illustrationblock}
\end{frame}

% FRAME 35
\begin{frame}[fragile]{12 | Complete a function pointer declaration -- Solution}
\begin{lstlisting}[style=solutioncpp]
int addition(int a, int b)
{
    return a + b;
}

int (*operation)(int, int) = addition;
int result = operation(1, 2);
\end{lstlisting}
\begin{columns}[T,onlytextwidth]
\begin{column}{0.48\textwidth}
\begin{definitionblock}{Compatible functions}
\texttt{subtraction}, \texttt{maximum}, or any other function with signature \texttt{int(int,int)} can be assigned.
\end{definitionblock}
\end{column}
\begin{column}{0.48\textwidth}
\begin{propositionblock}{Invariant call site}
The expression \texttt{operation(a,b)} stays unchanged while the selected behavior changes.
\end{propositionblock}
\end{column}
\end{columns}
\end{frame}

% FRAME 36
\begin{frame}{One CRR engine should not hard-code one payoff}
\begin{columns}[c,onlytextwidth]
\begin{column}{0.44\textwidth}
\centering
Model data and $q$\par
$\downarrow$\par
Terminal stock prices\par
$\downarrow$\par
\textbf{Selected payoff function}\par
$\downarrow$\par
Same backward induction
\end{column}
\begin{column}{0.52\textwidth}
\begin{definitionblock}{What varies}
At expiry, a call, put, or digital option applies a different function to the same stock value $S(N,i)$.
\end{definitionblock}
\begin{propositionblock}{What stays fixed}
\texttt{RiskNeutProb}, \texttt{CalculateAssetPrice}, the loops, discounting, and array storage do not change.
\end{propositionblock}
\end{column}
\end{columns}
\end{frame}

% FRAME 37
\begin{frame}[fragile]{13 | Add a payoff parameter to \texttt{PriceByCRR} -- Problem}
\problemmeta{Function pointer + dynamic array}{Complete the changing pieces}
\begin{lstlisting}[style=tinycpp]
double* PriceByCRR(double S0, double U, double D,
                   double R, int N, double K,
                   double (/* name */)(double, double))
{
    double q = RiskNeutProb(U, D, R);
    double* Price = new double[/* size */];
    for (int i = 0; i <= N; ++i)
    {
        double stock = CalculateAssetPrice(S0, U, D, N, i);
        Price[i] = /* call selected payoff */;
    }
    // backward induction follows
}
\end{lstlisting}
\end{frame}

% FRAME 38
\begin{frame}[fragile]{13 | The terminal row calls the selected payoff -- Solution A}
\begin{lstlisting}[style=tinycpp]
double* PriceByCRR(double S0, double U, double D,
                   double R, int N, double K,
                   double (*Payoff)(double, double))
{
    double q = RiskNeutProb(U, D, R);
    double* Price = new double[N + 1];
    for (int i = 0; i <= N; ++i)
    {
        double stock = CalculateAssetPrice(S0, U, D, N, i);
        Price[i] = Payoff(stock, K);
    }
\end{lstlisting}
\begin{propositionblock}{Single substitution point}
Only the terminal payoff evaluation depends on the option type.
\end{propositionblock}
\end{frame}

% FRAME 39
\begin{frame}[fragile]{13 | Backward induction remains unchanged -- Solution B}
\begin{lstlisting}[style=solutioncpp]
    for (int n = N - 1; n >= 0; --n)
    {
        for (int i = 0; i <= n; ++i)
        {
            Price[i] =
                (q * Price[i + 1]
                 + (1.0 - q) * Price[i]) / R;
        }
    }
    return Price;
}
\end{lstlisting}
\begin{definitionblock}{Ownership transfer}
The returned pointer owns an array created with \texttt{new[]}. The caller reads \texttt{Price[0]} and must eventually call \texttt{delete[]}.
\end{definitionblock}
\end{frame}

% FRAME 40
\begin{frame}[fragile]{14 | Complete call and put payoff functions -- Problem}
\problemmeta{Same signature, different rule}{Fill both branches}
\begin{lstlisting}
double CallPayoff(double z, double K)
{
    if (/* in the money */) return /* payoff */;
    return 0.0;
}

double PutPayoff(double z, double K)
{
    return (/* in the money */) ? /* payoff */ : 0.0;
}
\end{lstlisting}
\begin{definitionblock}{Acceptance cases at $K=100$}
$z=80$: call $0$, put $20$; \quad $z=120$: call $20$, put $0$.
\end{definitionblock}
\end{frame}

% FRAME 41
\begin{frame}[fragile]{14 | Complete call and put payoff functions -- Solution}
\begin{lstlisting}[style=solutioncpp]
double CallPayoff(double z, double K)
{
    if (z > K) return z - K;
    return 0.0;
}

double PutPayoff(double z, double K)
{
    return (z < K) ? (K - z) : 0.0;
}
\end{lstlisting}
\begin{propositionblock}{Why one pointer type accepts both}
Both functions return \texttt{double} and accept exactly two \texttt{double} parameters.
\end{propositionblock}
\end{frame}

% FRAME 42
\begin{frame}[fragile]{15 | Price a call and a put without leaking memory -- Problem}
\problemmeta{Lecture 3 driver}{Complete selection and cleanup}
\begin{lstlisting}[style=tinycpp]
double* optionPrice = nullptr;

optionPrice = PriceByCRR(S0, U, D, R, N, K,
                         /* call function */);
std::cout << optionPrice[0] << std::endl;
/* release array */;
optionPrice = nullptr;

optionPrice = PriceByCRR(S0, U, D, R, N, K,
                         /* put function */);
std::cout << optionPrice[0] << std::endl;
/* release array */;
optionPrice = nullptr;
\end{lstlisting}
\end{frame}

% FRAME 43
\begin{frame}[fragile]{15 | Price a call and a put without leaking memory -- Solution}
\begin{lstlisting}[style=tinycpp]
double* optionPrice = nullptr;

optionPrice = PriceByCRR(S0, U, D, R, N, K,
                         CallPayoff);
std::cout << optionPrice[0] << std::endl;
delete[] optionPrice;
optionPrice = nullptr;

optionPrice = PriceByCRR(S0, U, D, R, N, K,
                         PutPayoff);
std::cout << optionPrice[0] << std::endl;
delete[] optionPrice;
optionPrice = nullptr;
\end{lstlisting}
\begin{propositionblock}{Lecture 3 parameters}
For $S_0=106$, $U=1.15125$, $D=0.86862$, $R=1.00545$, $N=8$, $K=100$: call $=21.68$, put $=11.43$.
\end{propositionblock}
\end{frame}

% FRAME 44
\begin{frame}[fragile]{Overloading lets one name serve two input shapes}
\begin{lstlisting}
int GetInputData(double& S0, double& U,
                 double& D, double& R);

int GetInputData(int& N, double& K);

GetInputData(S0, U, D, R); // model overload
GetInputData(N, K);        // option overload
\end{lstlisting}
\begin{definitionblock}{How the compiler chooses}
The name is the same, but the number and types of arguments are different. The matching parameter list selects the overload.
\end{definitionblock}
\begin{illustrationblock}{Return type is not enough}
Two functions cannot be overloaded only by changing their return type.
\end{illustrationblock}
\end{frame}

% FRAME 45
\begin{frame}[fragile,shrink=4]{16 | Add digital options without changing CRR -- Problem}
\problemmeta{Lecture 3 extension}{Complete two new strategies}
\begin{lstlisting}
double DigitCallPayoff(double z, double K)
{
    return (/* condition */) ? 1.0 : 0.0;
}

double DigitPutPayoff(double z, double K)
{
    return (/* condition */) ? 1.0 : 0.0;
}

double* digitalCall = PriceByCRR(
    S0, U, D, R, N, K, /* function */);
\end{lstlisting}
\begin{definitionblock}{Boundary convention}
In the lecture, both digital payoffs return zero when $z=K$.
\end{definitionblock}
\end{frame}

% FRAME 46
\begin{frame}[fragile,shrink=3]{16 | Add digital options without changing CRR -- Solution}
\begin{lstlisting}[style=solutioncpp]
double DigitCallPayoff(double z, double K)
{
    return (z > K) ? 1.0 : 0.0;
}

double DigitPutPayoff(double z, double K)
{
    return (z < K) ? 1.0 : 0.0;
}

double* digitalCall = PriceByCRR(
    S0, U, D, R, N, K, DigitCallPayoff);
delete[] digitalCall;
\end{lstlisting}
\begin{propositionblock}{Open/closed design}
The pricer is open to a new payoff function but closed to edits in its backward-induction algorithm.
\end{propositionblock}
\end{frame}

% FRAME 47
\begin{frame}{The pointer story now forms one complete pricing workflow}
\begin{columns}[c,onlytextwidth]
\begin{column}{0.46\textwidth}
\centering
Address and dereference\par
$\downarrow$\par
Array and dynamic lifetime\par
$\downarrow$\par
Function pointer selects payoff\par
$\downarrow$\par
One CRR engine prices many claims
\end{column}
\begin{column}{0.50\textwidth}
\begin{propositionblock}{Three questions before running code}
\begin{enumerate}
\item What object or function does this pointer identify?
\item Is that target still alive and compatible?
\item Who owns any dynamic memory and releases it?
\end{enumerate}
\end{propositionblock}
\end{column}
\end{columns}
\end{frame}

% FRAME 48 -- reserve section
\begin{frame}{Brainteasers combine the Lecture 3 tools}
\begin{columns}[c,onlytextwidth]
\begin{column}{0.47\textwidth}
\begin{definitionblock}{For each challenge}
\begin{enumerate}
\item Trace the pointers by hand.
\item Complete the missing operations.
\item State the lifetime and bounds invariant.
\item Give one boundary test.
\end{enumerate}
\end{definitionblock}
\end{column}
\begin{column}{0.49\textwidth}
\begin{illustrationblock}{Difficulty ladder}
Reverse an array, compact values in place, solve a sorted pair sum, then select a payoff function dynamically.
\end{illustrationblock}
\medskip
These can be used as the final class challenge or as reserve practice.
\end{column}
\end{columns}
\end{frame}

% FRAME 49
\begin{frame}[fragile,shrink=6]{Brainteaser 01 | Reverse an array with two raw pointers -- Problem}
\problemmeta{In-place two pointers}{Complete movement and swap}
\begin{lstlisting}[style=tinycpp]
void Reverse(int* begin, int* end)
{
    while (begin < end)
    {
        /* move end from one-past-last to last */;
        if (begin >= end) break;
        int temp = /* left value */;
        *begin = /* right value */;
        *end = temp;
        /* move begin right */;
    }
}

int a[] = {1, 2, 3, 4, 5};
Reverse(a, a + 5);
\end{lstlisting}
\begin{definitionblock}{Target}
The array becomes \texttt{\{5,4,3,2,1\}} in $O(n)$ time and $O(1)$ extra space.
\end{definitionblock}
\end{frame}

% FRAME 50
\begin{frame}[fragile]{Brainteaser 01 | Reverse an array with two raw pointers -- Solution}
\begin{lstlisting}[style=solutioncpp]
void Reverse(int* begin, int* end)
{
    while (begin < end)
    {
        --end;
        if (begin >= end) break;
        int temp = *begin;
        *begin = *end;
        *end = temp;
        ++begin;
    }
}
\end{lstlisting}
\begin{propositionblock}{Invariant}
Elements outside the half-open interval \texttt{[begin,end)} are already in their final reversed positions.
\end{propositionblock}
\end{frame}

% FRAME 51
\begin{frame}[fragile]{Brainteaser 02 | Compact non-zero values in place -- Problem}
\problemmeta{Read and write pointers}{Return the new logical end}
\begin{lstlisting}[style=tinycpp]
int* CompactNonZero(int* begin, int* end)
{
    int* write = begin;
    for (int* read = begin; read != end; ++read)
    {
        if (/* keep this value */)
        {
            /* copy to output position */;
            /* advance output position */;
        }
    }
    return /* one-past-last kept value */;
}
\end{lstlisting}
\begin{definitionblock}{Example}
\texttt{\{0,5,0,-2,7\}} begins with \texttt{\{5,-2,7\}} after compaction; the returned pointer minus \texttt{begin} equals \texttt{3}.
\end{definitionblock}
\end{frame}

% FRAME 52
\begin{frame}[fragile]{Brainteaser 02 | Compact non-zero values in place -- Solution}
\begin{lstlisting}[style=solutioncpp]
int* CompactNonZero(int* begin, int* end)
{
    int* write = begin;
    for (int* read = begin; read != end; ++read)
    {
        if (*read != 0)
        {
            *write = *read;
            ++write;
        }
    }
    return write;
}
\end{lstlisting}
\begin{propositionblock}{Invariant}
\texttt{[begin,write)} contains exactly the non-zero values seen so far, in their original order.
\end{propositionblock}
\end{frame}

% FRAME 53
\begin{frame}[fragile]{Brainteaser 03 | Find a target sum in a sorted array -- Problem}
\problemmeta{Const pointers + linear scan}{Move the correct side}
\begin{lstlisting}[style=tinycpp]
bool HasPairSum(const int* values, int n, int target)
{
    const int* left = values;
    const int* right = values + n - 1;
    while (left < right)
    {
        int sum = *left + *right;
        if (sum == target) return true;
        if (sum < target) /* move left */;
        else              /* move right */;
    }
    return false;
}
\end{lstlisting}
\begin{definitionblock}{Example}
\texttt{\{1,2,4,7,11\}}, target \texttt{9} returns \texttt{true} from $2+7$.
\end{definitionblock}
\end{frame}

% FRAME 54
\begin{frame}[fragile]{Brainteaser 03 | Find a target sum in a sorted array -- Solution}
\begin{lstlisting}[style=solutioncpp]
if (sum == target) return true;
if (sum < target) ++left;
else              --right;
\end{lstlisting}
\begin{columns}[T,onlytextwidth]
\begin{column}{0.48\textwidth}
\begin{definitionblock}{Why moving left works}
The array is sorted. If the sum is too small, decreasing the right value cannot help; increase the left value.
\end{definitionblock}
\end{column}
\begin{column}{0.48\textwidth}
\begin{propositionblock}{Complexity}
Each pointer moves inward at most $n-1$ times: $O(n)$ time and $O(1)$ extra space.
\end{propositionblock}
\end{column}
\end{columns}
\end{frame}

% FRAME 55
\begin{frame}[fragile]{Brainteaser 04 | Return the selected payoff function -- Problem}
\problemmeta{Function pointer as a return value}{Complete three branches}
\begin{lstlisting}[style=tinycpp]
using Payoff = double (*)(double, double);

Payoff SelectPayoff(char code)
{
    if (code == 'C') return /* call */;
    if (code == 'P') return /* put */;
    return /* no valid function */;
}

Payoff payoff = SelectPayoff('P');
if (payoff != nullptr)
    std::cout << payoff(80.0, 100.0);
\end{lstlisting}
\end{frame}

% FRAME 56
\begin{frame}[fragile]{Brainteaser 04 | Return the selected payoff function -- Solution}
\begin{lstlisting}[style=solutioncpp]
using Payoff = double (*)(double, double);

Payoff SelectPayoff(char code)
{
    if (code == 'C') return CallPayoff;
    if (code == 'P') return PutPayoff;
    return nullptr;
}
\end{lstlisting}
\begin{definitionblock}{Safe call site}
Check the returned pointer before invoking it. With code \texttt{'P'}, \texttt{payoff(80,100)} returns \texttt{20}.
\end{definitionblock}
\begin{illustrationblock}{Extension}
Add \texttt{'D'} for a digital call without changing the calling code.
\end{illustrationblock}
\end{frame}

% FRAME 57
\begin{frame}{The final test is to explain each star and each delete}
\begin{columns}[c,onlytextwidth]
\begin{column}{0.47\textwidth}
\begin{definitionblock}{Exit ticket}
\begin{enumerate}
\item What does \texttt{*} mean here?
\item Is the target alive and in bounds?
\item Does \texttt{const} protect the pointer or the value?
\item Who owns the allocation?
\end{enumerate}
\end{definitionblock}
\end{column}
\begin{column}{0.49\textwidth}
\begin{propositionblock}{Transfer to CRR}
The payoff function pointer changes the terminal condition. It does not change the tree mechanics, risk-neutral probability, or backward induction.
\end{propositionblock}
\medskip
If you can answer the four questions before running the program, you can debug pointer code instead of guessing at addresses.
\end{column}
\end{columns}
\end{frame}

\end{document}
